\section{Horizontal leaves and their finite monodromy}\label{sec:horizontal}
We construct compact symplectic six-manifolds with integral first Chern class zero for which an entire face of the universal cycle series can be computed in every covering degree. The geometric factors have unbounded genera in a single compact example. In the corresponding local models, passing to all finite covers of the base recovers exactly the monodromy of any finite collection of sheets, and characterizes the integer multiplicities by subgroup congruences. We use the smooth universal cycle group \(S\) of the geometric product formula of \cite{WueProduct}. Write \(E_h=\Phi(A_h)\), where \(A_h\) denotes the balanced antidiagonal series. Its coefficient of order one is \(\Phi(x_{h,1})\ne0\) \cite{WueProduct}.

\subsection{The local model in every genus}

Let \(C=\Gamma\backslash\mathbb H\) be a closed Riemann surface of genus \(g\ge2\). Choose a lift of its uniformizing representation to \(\mathrm{SL}_2(\mathbb R)\). Such lifts correspond to spin structures: the obstruction to lifting the oriented tangent circle bundle to its spin double cover is the Euler characteristic modulo two, which vanishes. Let \(V\) be its flat real plane bundle in the standard representation. Choose the sign of its invariant area form so that the oriented plane has Euler number \(g-1\). Indeed the projective-line circle bundle of \(V\) is the ideal-boundary bundle of the hyperbolic surface. It is identified with the unit tangent circle bundle by sending a direction to the endpoint of its geodesic. Passing from oriented unit vectors in \(V\) to their unoriented lines doubles the Euler class up to orientation. Thus its absolute Euler number is \(g-1\). Choose a compatible complex structure on \(V\) varying smoothly over \(C\). On the flat bundle \(V\oplus V\) take the product vertical complex structure and take the base complex structure on the horizontal distribution. The sum of the two constant vertical area forms in flat coordinates is a global closed form \(\alpha\). It is nonnegative on every complex tangent line and vanishes precisely on horizontal lines. In this vector bundle \(\alpha\) is exact by the radial Liouville primitive. Every compact nonconstant holomorphic curve is consequently horizontal.

This argument is unchanged on a finite cover of the base, where each of the two normal complex lines has degree $g(T)-1$ on the covering curve $T$, and their direct sum has degree $2g(T)-2$. For a branched holomorphic cover \(f:D\to C\), the graph equation for a normal section is the pullback of the real Cauchy--Riemann operator defined by the flat connection and the chosen fiber complex structure. Its kernel consists of parallel sections: the same exact area form on the pullback vanishes on the entire holomorphic graph. The image of \(\pi_1(D)\) has finite index, since \(f\) lifts to the associated connected covering surface, where its image is both compact and open. Any finite-index subgroup of \(\Gamma\) contains a nonzero power of a hyperbolic element, which fixes no nonzero vector in the standard linear representation. Hence every pullback operator is injective. The two summands have the \emph{same} real operator, including its orientation data. Furthermore a finite orbit in the linear fiber is just the origin, by applying a power of the same hyperbolic element. Thus the core is the sole irreducible curve in this model. The doubled elementary-neighborhood theorem of \cite{WueProduct} identifies the actual characteristic insertion of its \(m\)-fold core with \([q^m]E_g\) for every \(m\). This application uses the full universal elementary-neighborhood result.

For any finite \(\Gamma\)-set \(\Omega\) form the flat suspension with fiber \(\Omega\times(\mathbb R^2\oplus\mathbb R^2)\), acting by permutation on the first factor and by the doubled representation on the second. This is a union of vector-bundle neighborhoods over the connected covers associated with the orbits. The reduced curves in it are precisely the zero sections of those covers. If \(O\) is an orbit of cardinality \(d_O\), its curve \(D_O\) has genus \(1+d_O(g-1)\). Retaining all actual homology labels, the support series is
\[\prod_{O\in\Gamma\backslash\Omega}
                  E_{1+d_O(g-1)}(Q^{[D_O]}).\]
Indeed a cycle in any fixed total degree is a choice of nonnegative multiples of the disjoint zero sections. It gives the product of their open compact characteristic insertions by excision and disjoint union. No finiteness in source genus is used to identify this cycle space. In this class of repeated normal operators it is the geometry of the single zero section and all its multiples that fixes the factor attached to a transitive support.

\subsection{Detecting the support under finite base change}

Retain only the projection degree in the displayed product. The resulting series determines the multiset of cardinalities of the orbits. For two different such multisets, take the smallest cardinality \(d\) at which the counts differ. The difference of the coefficients of \(q^d\) is that nonzero integer difference times \(\Phi(x_{1+d(g-1),1})\). Contributions from smaller orbits agree and the class just displayed is nonzero. Conversely equal cardinality multisets plainly give the same projected series.

For a finite-index subgroup \(K\le\Gamma\), pull the neighborhood back to the corresponding connected unramified cover \(C_K\). Its sheet set is the restriction of the \(\Gamma\)-action to \(K\). Its coefficient of projection degree one over \(C_K\) is
\[|\Omega^K|\,\Phi(x_{g(C_K),1}).\]
Thus the collection of these coefficients for all finite covers determines the entire isomorphism class of \(\Omega\) as a finite \(\Gamma\)-set. To see this directly for two candidates take a common finite quotient \(G\) through which both actions factor. Every subgroup \(H\le G\) is tested by its inverse image in \(\Gamma\). If \(m_H=|\Omega^H|\), Möbius inversion in the finite subgroup lattice gives the number of points with stabilizer \emph{equal} to \(H\) as
\[t_H=\sum_{H\le J\le G}\mu(H,J)m_J .\]
These points occur in groups of \([N_G(H):H]\) in each orbit isomorphic to \(G/H\). For conjugacy-invariant integer data \(m_H\), the conditions \(t_H\ge0\) and \([N_G(H):H]\mid t_H\) are also sufficient. Take \(t_H/[N_G(H):H]\) disjoint copies of \(G/H\) for a representative of each conjugacy class and invert the formula. Dropping the inequalities gives the Grothendieck group, the additive group of the Burnside ring of finite \(G\)-sets. Under restriction these constructions are ordinary base change of finite covers; induction from \(K\) to \(\Gamma\) is composition of the corresponding covers. Fiber product of covers gives multiplication of the marks. Every finite set used in this inversion is geometrically realized by the preceding suspension, including when the union of neighborhoods is disconnected. The integrality conditions here come from stabilizers of actual sheets.

One series over \(C\) alone does lose information. Take two distinct index-two subgroups \(K_1,K_2\) of the surface group. They are normal and are not conjugate. The two transitive two-sheet models give identical projected series \(E_{2g-1}(q^2)\). On the cover for \(K_1\) the first has two degree-one leaves while the second has none. This distinguishes geometric support data even when every coefficient in the original projected series agrees.

\subsection{A compact symplectic Calabi–Yau application}

Take a rational quaternion division algebra split at the real place, a maximal order \(\Lambda\) and a torsion-free principal congruence subgroup \(\Gamma\) of the norm-one units, of level divisible by an integer at least three chosen to eliminate torsion. The arithmetic Fuchsian quotient \(C=\Gamma\backslash\mathbb H\) is compact. Under a real splitting, the four-dimensional real algebra is \(M_2(\mathbb R)\) and left multiplication is the direct sum of the standard representations on its two columns. The order is a rank-four lattice preserved by this action. Put
\[X=(\mathbb H\times(M_2(\mathbb R)/\Lambda))/\Gamma .\]
The smooth torus bundle has a zero section \(C_0\). Give its horizontal and vertical distributions the almost-complex structures above. The form \(\alpha\) descends to the torus bundle. If \(\omega_C\) is an area form on \(C\), then \(\alpha+\pi^*\omega_C\) is a compatible symplectic form. The first Chern class vanishes \emph{integrally} since the complex tangent splitting is the pullback of \(TC\oplus V\oplus V\) and its first Chern number on \(C\) is \(2-2g+2(g-1)=0\). These statements use the chosen almost-complex structure; the resulting class formula is in \(S\). The classical quaternionic abelian surface construction supplies the lattice and the compact hyperbolic base \cite{voight}.

A pseudoholomorphic irreducible curve with \(\alpha\)-area zero is horizontal by pointwise nonnegativity on its normalization. It lies in one horizontal leaf and covers that leaf, so its reduced image is a compact leaf of a finite orbit. One can see the last implication by lifting base paths in the leaf: the normalization maps nonconstantly into its covering surface, and is open and compact there. Conversely each finite orbit gives a compact embedded holomorphic leaf. The finite-orbit points in the torus are precisely its torsion subgroup. An orbit of size at most \(n\) is fixed by \(\gamma^{n!}\) for a single chosen hyperbolic \(\gamma\in\Gamma\). The matrix \(\gamma^{n!}-1\) is invertible on the real four-space and integral in \(\Lambda\), so its kernel on the torus is a finite torsion group. Every torsion point lies in a finite invariant torsion subgroup of the torus. In particular only finitely many leaves occur in bounded projection degree.

A leaf associated with a point \(p\) of stabilizer \(\Gamma_p\) has an embedded neighborhood of the zero section in the local model over \(C_p\). The map from its flat vector bundle to \(X\) sends \(v\) to \(p+v\) in the torus. It is a local diffeomorphism along the compact section, embeds that section, and restricts to an embedding on a small tubular neighborhood by compactness. It preserves the horizontal and vertical complex structures. It follows that for every integral homology class \(\beta\) with \(\langle\alpha,\beta\rangle=0\) the universal insertion of the entire cycle space in that class has generating series
\[\sum_{\langle\alpha,\beta\rangle=0} [Z(X,J;\beta)]\,Q^\beta
  = \prod_{O\in\Gamma\backslash(M_2(\mathbb R)/\Lambda)_{\rm tors}}
        E_{1+|O|(g-1)}(Q^{[D_O]}).\]
On the left the bracket denotes the image in \(S\) of the unit of that compact open-and-closed class locus, including the empty cycle. Homology labels of integral cycles are locally constant in the cycle topology, as in \cite{WueProduct}; alternatively each contributing cycle has disjoint tubular neighborhoods with only its prescribed leaf and multiples. For the latter argument sufficiently small perturbations of the cycle keep the same component degrees in those tubes. Compactness in each class here also follows directly from the finite-orbit bound. Completion by projection degree on this face makes every coefficient finite and does not depend on the choice of area form on \(C\). This face has infinitely many leaves with unbounded positive genus.

\subsection{Homology under multiplication and an infinite arithmetic family}

Let \(P_n:X\to X\) be multiplication by \(n\) in the torus fibers. For any orbit in \(X[n]\) its leaf of degree \(d\) maps with degree \(d\) to \(C_0\). In particular the actual class \([D_O]\) satisfies \(P_{n*}[D_O]=d[C_0]\). One also has the integral bound
\[n^3\bigl([D_O]-d[C_0]\bigr)=0.\]
Indeed \(\delta=[D_O]-d[C_0]\) projects to zero on \(H_2(C;\mathbb Z)\), hence is in filtration step \(F_1\) of the homological Serre spectral sequence of the torus bundle. The edge quotient \(F_2/F_1\) injects into \(H_2(C;\mathbb Z)\); the zero section splits the projection. On \(F_1/F_0\) multiplication by \(n\) acts by \(n\), and on \(F_0\) by \(n^2\), since those are subquotients of the terms with torus fiber degrees one and two. Thus \((P_{n*}-n^2)(P_{n*}-n)\delta=0\). Since \(P_{n*}\delta=0\), the bound follows. It does not assert equality of torsion-distinct labels at the exceptional primes.

At any prime \(p\) away from the quaternion discriminant and the congruence level, reduction of \(\Gamma\) acts on \(X[p]\simeq M_2(\mathbb F_p)\) by the full group \(\mathrm{SL}_2(\mathbb F_p)\) on the left. This follows by strong approximation for the norm-one group at the split real place, with the prescribed congruence condition away from \(p\) and the chosen residue at \(p\). The maximal local order is \(M_2(\mathbb Z_p)\) and the reduction of its determinant-one group is surjective \cite{voight}.

The zero matrix gives the section. Rank-one matrices have \(p+1\) orbits, indexed by their row line, each of size \(d_1=p^2-1\). After fixing one generator of that row line the column is any nonzero vector, on which \(\mathrm{SL}_2\) acts transitively. Invertible matrices have \(p-1\) orbits, indexed by determinant, each of size \(d_2=p(p^2-1)\). The sizes add to \(p^4\) after including orbit multiplicities and the zero matrix. The union of these leaves is exactly the preimage of the zero section under \(P_p\). It is compact and has a product of the elementary contributions of its disjoint leaves in every multiple. For every such prime the genuine homology classes group by degree. Here is the integral check. Set \(U=\mathbb F_p^2\). The reduced fiber lattice is \(U\oplus U\) for the full group \(\mathrm{SL}_2(\mathbb F_p)\), so it has no invariants. Thus \(H^1(C;\Lambda)\) has no \(p\)-torsion by the long exact sequence for multiplication by \(p\), and surface Poincaré duality gives the same assertion for \(H_1(C;\Lambda)\). The coinvariants \((\bigwedge^2\Lambda)_\Gamma\) have dimension three after reduction modulo \(p\). Indeed the wedge of the direct sum decomposes as two determinant lines and \(U\otimes U\); the latter has one-dimensional coinvariants since \(U\) is symplectically self-dual and the centralizer of the two elementary transvections consists of scalars. This works also in characteristic two. Over the reals the coinvariants have dimension at least three from the two invariant column area forms and the invariant alternating form pairing the two columns by their determinant. Their integral rank is therefore three and there is no \(p\)-torsion. In the homological Serre sequence the only possible incoming differential to this fiber-degree-two term is from \(H_2(C;\Lambda)=\Lambda^\Gamma=0\) by Poincaré duality. The \((1,1)\) term is untouched and the zero section splits the base term. Hence \(H_2(X;\mathbb Z)\) has no \(p\)-torsion at every good prime. Applying the preceding \(n^3\) bound for \(n=p\) proves the claim. With \(q=Q^{[C_0]}\), the series supported on that preimage, in the actual integral-class completion, is then
\[E_g(q)\,
     E_{1+d_1(g-1)}(q^{d_1})^{p+1}\,
     E_{1+d_2(g-1)}(q^{d_2})^{p-1}.\]
At every good prime the same expression holds after retaining just the projection degree. For orders divisible by a prime at the level or discriminant the individual orbit product retains all integral labels. The formula concerns the indicated finite preimage inside the horizontal face; other torsion orders also contribute to the entire face.

There is also a formula for the entire part of the horizontal face supported on points of order prime to the level and discriminant. At a good prime and exact order \(p^k\), set \(R=\mathbb Z/p^k\) and \(u_k=p^{2k}-p^{2k-2}\). The exact-order points are the matrices over \(R\) with a unit entry. Put \(t=\min(v_p(\det),k)\), with \(t=k\) for zero determinant. Every orbit with this value has degree \(u_k p^{k-t}\). The number of these orbits is
\[\begin{cases}
  \varphi(p^k), &t=0,\\
  \varphi(p^{k-t})(p^t+p^{t-1}), &0<t<k,\\
  p^k+p^{k-1}, &t=k .
 \end{cases}\]
Indeed, if the first column is primitive, normalize it to the first standard basis vector by \(\mathrm{SL}_2(R)\). There are \(u_k\) primitive columns. The other column becomes \((a,\det)\), and the stabilizer adds multiples of \(\det\) to \(a\). This gives \(p^t\) choices of its residue and \(p^{k-t}\) points per full fiber of that orbit. If the first column is nonprimitive, normalize the second one instead. Now \(p\mid a\) and \(p\mid\det\), giving \(p^{t-1}\) more residues when \(t>0\). Counting determinants proves the formula. For any composite good order \(n\), strong approximation and the Chinese remainder theorem give the product of the prime-power actions on the product of the exact-order point sets. Choose one row of this formula at each prime dividing \(n\) and multiply the orbit degrees and the orbit counts separately. An orbit of resulting degree \(d\) contributes \(E_{1+d(g-1)}(q^d)\); its exponent is the resulting count. The true label is always \(d[C_0]\) by the integral calculation at the primes dividing \(n\). The product over exact orders is locally finite by the single-hyperbolic-element argument above.

The construction identifies the geometric integral structure and the information lost on forgetting sheets in this class of higher-genus suspensions. Its local normal operators are doubled and injective on all covers. A higher-genus wall law at failure of that injectivity, and a complex-algebra lift of general compact symplectic face identities, require different arguments.
